
\subsubsection{Root Causes}

\textbf{Cause 1: Insufficient Hutchinson Probes}

Bekas et al. (2007) establish O(1/ε²) sample complexity for ε-accuracy in trace estimation. For 768-dimensional embeddings and target CV <15\%, this implies ~100+ probes. Our implementation used 10 probes, likely producing high-variance estimates. The degenerate zero measurements suggest either extreme variance or gradient detachment.

\textbf{Cause 2: Power Iteration Non-Convergence}

Five iterations may be insufficient for spectral norm convergence in residual-corrected Jacobians through attention and LayerNorm operations. The zero spectral norms across all layers and checkpoints suggest non-convergence or numerical underflow in the computation graph.

\textbf{Cause 3: Autodiff Gradient Pathologies}

Computing layer-wise Jacobians requires backpropagation through attention mechanisms, MLPs, and LayerNorm. Differentiating through the entire spectral estimation procedure (Hutchinson + power iteration) creates deep computation graphs. The regularization loss magnitude (-10^10) exceeding CLM loss (~10) by nine orders of magnitude indicates gradient explosion. The negative sign violates the mathematical constraint that stable rank is non-negative (sr ≥ 0 by construction).

\subsubsection{Measurement-Control Gap}

Post-hoc SVD on saved activations is numerically stable because it operates on fixed tensors without gradient flow. Our approach computes stable rank during forward pass with gradients enabled, feeding it into the loss, and backpropagating through the entire measurement pipeline. This tight coupling means measurement errors directly corrupt the training signal.

\subsubsection{Comparison to Successful Methods}

Spectral normalization (Miyato et al., 2018) succeeds because it operates on static weight matrices, not dynamic Jacobians. LoRA and ARD-LoRA succeed because they optimize rank on small adapter matrices with explicit low-rank structure, not via stochastic spectral estimation on full layers.
